\section{Assignment questions awaiting student responses}\label{sec:questions}
The assignment requires four experiment-supported answers and two literature-review answers. No response is supplied here. Complete each placeholder using the indicated evidence, and cite the exact analysis/output added to the final submission.

\subsection{Mandatory question 1: leakage and preprocessing}
\textbf{[ANSWER TO BE COMPLETED BY STUDENT]}\par
Required evidence: run a small controlled development-only comparison showing imputation or SMOTE applied before splitting versus fitted/applied inside each training fold. Report fold scores, explain leakage direction, and identify which estimate is trustworthy. Do not use the held-out test set for this experiment.

\subsection{Mandatory question 2: probabilities after 1:1 undersampling}
In five-fold stratified development cross-validation, the observed mortality prevalence across the 1,360 outer-fold predictions was 0.1596. A Stage~2 regularized logistic regression was trained after 1:1 undersampling. For each outer fold, a separate stratified 20\% subset of the outer training data retained its natural prevalence; the base model, fit on the disjoint remaining subset after undersampling, generated scores for this calibration subset. Platt scaling (Section~11.3) was fitted to those held-out scores and labels, then applied to the untouched outer validation fold. No validation labels were used for fitting. The pooled mean predicted mortality probabilities were 0.3411 raw, 0.1768 after the distinct odds-based prior correction, and 0.1663 after Platt calibration. Brier score was 0.1520, 0.0964, and 0.0938, respectively; log loss was 0.4802, 0.3246, and 0.3103; ten-bin quantile ECE was 0.1831, 0.0436, and 0.0123. Platt had the best pooled probability scores and a calibration slope of 1.0116/intercept of -0.0563, though its AUROC (0.8534) was slightly lower than raw/prior correction (about 0.8590). Its mean probability (0.1663) remained slightly above prevalence. The outputs, reliability plot, and run manifest are in \texttt{undersampling\_calibration\_outputs/}; the earlier prior-correction-only experiment is preserved separately. No held-out records were used.

\subsection{Mandatory question 3: asymmetric decision cost}
Using the frozen T06 CatBoost model's Platt-calibrated five-fold outer development OOF probabilities (1,360 patients; observed death prevalence 0.1596), I classify a patient as high risk when $p(\mathrm{death})\ge t$ and minimize empirical cost $10FN+FP$. The threshold grid included zero, every distinct OOF probability, and an all-negative cutoff just above the maximum; ties select the highest threshold. The minimum-cost threshold was $t=0.101369$, with TP=197, FP=305, TN=838, FN=20, and total cost 505. Thus 502 patients were alerted: 36.91 alerts per 100 (369.12 per 1,000). There were 197 deaths identified: 14.49 per 100 (144.85 per 1,000), with 20 deaths missed. Alert rate is alerts divided by all patients, whereas sensitivity is deaths identified divided by actual deaths (0.9078); deaths identified per 100 patients is TP divided by all patients. At $t=0.5$, cost was 1,079; at the existing frozen threshold 0.153336, cost was 636. The selected threshold is an empirical development operating point evaluated on the same OOF records used to optimize it; its apparent cost may be optimistic, so it does not replace the frozen T06 threshold or constitute a final test evaluation. Results and reproducibility details are in \texttt{mi\_cost\_sensitive\_threshold\_outputs\_v2/}.

\subsection{Mandatory question 4: prediction versus causality}
\textbf{[ANSWER TO BE COMPLETED BY STUDENT]}\par
Required evidence: use the guide's heparin example and identify an important feature in the fitted model, then distinguish predictive association from causal effect. Explain confounding, treatment timing, and what causal design or data would be needed. The SHAP results in this report do not establish causality.

\subsection{Literature-review question 1: student-selected guide article}
\textbf{[ANSWER TO BE COMPLETED BY STUDENT]}\par
Select an article and its associated first literature-review prompt from the guide; summarize the study question, cohort, method, principal result, and limitations, with a citation. Article and prompt selection are pending.

\subsection{Literature-review question 2: student-selected guide article}
\textbf{[ANSWER TO BE COMPLETED BY STUDENT]}\par
Select an article and its associated second literature-review prompt from the guide; answer using the paper's evidence and a comparison with this project's task where appropriate. Article and prompt selection are pending; no literature questions have been selected on the student's behalf.
